\documentclass[10pt,a4paper]{amsart}
\usepackage{amsmath,amssymb,amsthm,graphicx}
% Retained minimal preamble for the Stacks Project chapter more-morphisms.tex (More on Morphisms).
% Verbatim excerpt of the chapter preamble: Xy-pic setup (native preamble lines 14-18) and the native
% theorem declarations with the mathematical macros (native preamble lines 154-217). Copied mechanically;
% nothing added, removed or redefined. The wrapper class and the amsmath/amssymb/amsthm/graphicx base
% packages are supplied by the assembler, so the class/externaldocument part of the native preamble is omitted.
% For commutative diagrams we use Xy-pic
\usepackage[all]{xy}

% We use 2cell for 2-commutative diagrams.
\xyoption{2cell}
% Theorem environments.
%
\theoremstyle{plain}
\newtheorem{theorem}[subsection]{Theorem}
\newtheorem{proposition}[subsection]{Proposition}
\newtheorem{lemma}[subsection]{Lemma}

\theoremstyle{definition}
\newtheorem{definition}[subsection]{Definition}
\newtheorem{example}[subsection]{Example}
\newtheorem{exercise}[subsection]{Exercise}
\newtheorem{situation}[subsection]{Situation}

\theoremstyle{remark}
\newtheorem{remark}[subsection]{Remark}
\newtheorem{remarks}[subsection]{Remarks}

\numberwithin{equation}{subsection}

% Macros
%
\def\lim{\mathop{\mathrm{lim}}\nolimits}
\def\colim{\mathop{\mathrm{colim}}\nolimits}
\def\Spec{\mathop{\mathrm{Spec}}}
\def\Hom{\mathop{\mathrm{Hom}}\nolimits}
\def\Ext{\mathop{\mathrm{Ext}}\nolimits}
\def\SheafHom{\mathop{\mathcal{H}\!\mathit{om}}\nolimits}
\def\SheafExt{\mathop{\mathcal{E}\!\mathit{xt}}\nolimits}
\def\Sch{\mathit{Sch}}
\def\Mor{\mathop{\mathrm{Mor}}\nolimits}
\def\Ob{\mathop{\mathrm{Ob}}\nolimits}
\def\Sh{\mathop{\mathit{Sh}}\nolimits}
\def\NL{\mathop{N\!L}\nolimits}
\def\CH{\mathop{\mathrm{CH}}\nolimits}
\def\proetale{{pro\text{-}\acute{e}tale}}
\def\etale{{\acute{e}tale}}
\def\QCoh{\mathit{QCoh}}
\def\Ker{\mathop{\mathrm{Ker}}}
\def\Im{\mathop{\mathrm{Im}}}
\def\Coker{\mathop{\mathrm{Coker}}}
\def\Coim{\mathop{\mathrm{Coim}}}

% Boxtimes
%
\DeclareMathSymbol{\boxtimes}{\mathbin}{AMSa}{"02}

%
% Macros for moduli stacks/spaces
%
\def\QCohstack{\mathcal{QC}\!\mathit{oh}}
\def\Cohstack{\mathcal{C}\!\mathit{oh}}
\def\Spacesstack{\mathcal{S}\!\mathit{paces}}
\def\Quotfunctor{\mathrm{Quot}}
\def\Hilbfunctor{\mathrm{Hilb}}
\def\Curvesstack{\mathcal{C}\!\mathit{urves}}
\def\Polarizedstack{\mathcal{P}\!\mathit{olarized}}
\def\Complexesstack{\mathcal{C}\!\mathit{omplexes}}
% \Pic is the operator that assigns to X its picard group, usage \Pic(X)
% \Picardstack_{X/B} denotes the Picard stack of X over B
% \Picardfunctor_{X/B} denotes the Picard functor of X over B
\def\Pic{\mathop{\mathrm{Pic}}\nolimits}
\def\Picardstack{\mathcal{P}\!\mathit{ic}}
\def\Picardfunctor{\mathrm{Pic}}
\def\Deformationcategory{\mathcal{D}\!\mathit{ef}}
\newenvironment{reference}{}{}
\newenvironment{slogan}{}{}
\newenvironment{history}{}{}
\usepackage{cleveref}
\crefname{originallemma}{original lemma}{original lemmas}
\Crefname{originallemma}{Original lemma}{Original lemmas}
\makeatletter
\newlabel{schemes-lemma-points-specialize}{{01J8}{source}{}{sourcepointer-1}{}}
\newlabel{schemes-lemma-points-specialize@cref}{{[originallemma][0][]01J8}{[1][0][]source}{}{}{}}
\makeatother
\title{Realising a generic point of a fibre component over a valuation ring base change}
\author{The Stacks Project Authors}
\date{stacks-project:more-morphisms.tex:local-snapshot-20261011; GNU FDL 1.2 or later}
\begin{document}
\maketitle
\noindent\textit{Editorial note.} Statement span is lines 5093-5125 of sources/stacks-more-morphisms/pinned-original-source.tex; proof span is lines 5127-5164, the complete author proof (36 body lines).\par\medskip
\noindent\textit{Editorial note.} Navigation stubs: 1 cited-but-not-included labels are supplied as non-load-bearing pointers with their pinned tags where the snapshot tag index has them.\par\medskip
\noindent\textit{Editorial note.} extra\_packages is empty on purpose: the retained preamble already loads xy with options and a second declaration clashes in the assembler wrapper.\par\medskip
\noindent\textit{Editorial note.} Text is retained verbatim, including the author's own typos and placeholders.\par\medskip
\noindent\textit{Editorial note.} The author's statement contains a reference environment citing EGA, so the card carries a bibliography fragment. That fragment reproduces, field for field and without editing, the @ARTICLE\{EGA\} entry from the Stacks Project snapshot bibliography shipped with this corpus at focused-seeds/AGP/sources/stacks-snapshot-20261011/my.bib (same upstream Stacks Project snapshot). No entry text was altered and no citation was reconstructed from memory.\par\medskip
\noindent\textit{Allowed standard prerequisites (editorial).}\begin{itemize}
\item Explicit prerequisite boundary. Spectra of valuation rings, closures viewed as integral subschemes, and generic points of components of a fibre; the construction is retained.
\end{itemize}\medskip
\ifdefined\hypertarget\hypertarget{sourcepointer-1}{}\fi
\noindent\textit{Source reference.} Original lemma 01J8: sources/stacks-more-morphisms/pinned-original-source.tex, source lines 5164--5164. Cited by the retained statement or proof but not part of the collected spans; supplied only as a navigation pointer. The argument does not depend on its content.\par\medskip
\begin{lemma}
\label{lemma-make-base-change}
\begin{reference}
This can be found in the proof of
\cite[IV Proposition 12.1.1.5]{EGA}
\end{reference}
Let $f : X \to S$ be a morphism of schemes. Let $y \in X$ be a point
with image $t \in S$. Denote $Y \subset X$ the closure of $\{y\}$
viewed as an integral closed subscheme of $X$. Let $s \in S$ and let
$x \in Y_s$ be a generic point of an irreducible component of $Y_s$.
There exists a cartesian diagram
$$
\xymatrix{
X' \ar[r]_{g'} \ar[d]_{f'} & X \ar[d]^f \\
S' \ar[r]^g & S
}
$$
with the following properties:
\begin{enumerate}
\item $S'$ is the spectrum of a valuation ring
with generic point $t'$ and closed point $s'$,
\item $g(t') = t$ and $g(s') = s$,
\item there exists a point $y' \in X'_{t'}$ which is
a generic point of an irreducible component of
$(S' \times_S Y)_{t'} = Y_t \times_t t'$
and satisfies $g'(y') = y$,
\item denoting $Y' \subset X'$ the closure of $\{y'\}$
viewed as an integral closed subscheme of $X'$
there exists a point $x' \in Y'_{s'}$ which is a generic
point of an irreducible component of $Y'_{s'}$
with $g'(x') = x$.
\end{enumerate}
\end{lemma}

\begin{proof}
We choose a valuation ring $R$, we set $S' = \Spec(R)$ with
generic point $t'$ and closed point $s'$, and we choose a morphism
$h : S' \to X$ with $h(t') = y$ and $h(s') = x$.
See Schemes, Lemma \ref{schemes-lemma-points-specialize}.
Set $g = f \circ h$ so that $g(t') = t$ and $g(s') = s$.
Consider the base change
$$
\xymatrix{
X' \ar[r]_{g'} \ar[d] & X \ar[d] \\
S' \ar@/^1em/[u]^\sigma \ar[r]^-g & S
}
$$
We obtain a section $\sigma$ of the base change such that
$h = g' \circ \sigma$.

\medskip\noindent
Of course $\sigma$ factors through the base change $S' \times_S Y$ of $Y$
as $h$ factors through $Y$. Let $y' \in X'_{t'} \subset X'$ be the
generic point of an irreducible component of the fibre
$$
(S' \times_S Y)_{t'} = Y_t \times_t t'
$$
containing the point $\sigma(t')$, i.e., such that
$y' \leadsto \sigma(t')$.
Since $g'(y') \in Y_t$ and $g(y') \leadsto g(\sigma(t')) = y$
we find that $g'(y') = y$ because $y$ is the generic point of
the fibre $Y_t$.
Denote $Y' \subset X'$ the closure of $\{y'\}$ in $X'$ viewed
as an integral closed subscheme.
Then $\sigma$ factors through $Y'$ as $\sigma(t') \in Y'$.
Choose a generic point $x' \in Y'_{s'}$ of an irreducible
component of $Y'_{s'}$ which contains $\sigma(s')$, i.e., we
get $x' \leadsto \sigma(s')$ and hence $g'(x') \leadsto g'(\sigma(s')) = x$.
Again as $x$ is a generic point of an irreducible component of $Y_s$
by assumption
and as $g'(Y') \subset Y$ we conclude that $g'(x') = x$.
\end{proof}

\begin{thebibliography}{1}
\bibitem{EGA}
Dieudonn\'{e}, Jean; Grothendieck, Alexander.
\newblock {\em \'{E}l\'ements de g\'eom\'etrie alg\'ebrique}.
\newblock Inst. Hautes \'Etudes Sci. Publ. Math., volumes 4, 8, 11, 17, 20, 24, 28, 32, 1961--1967.
\end{thebibliography}
\end{document}
